\cotefile{83}
\scanpage{17}
\mspage{7\handwritten{5} \deleted{\unsure{bis}}}
% archive sheet number: 16

J'ai repris \inserted{\handwritten{tantôt}} l'expression ``psychodrame'', terme assez à la mode je \deleted{\unsure{croi}}
crois\handwritten{,} et qui ne manque pas de force\handwritten{.} \deleted{évocatrice.} Et il est vrai que scénario
et mise en scène de beaucoup de rêves\deleted{, \unsure{ceux surtout qui virent dans le cauchemar,}}
ne ménagent pas les sombres couleurs, aux tonalités \deleted{\unsure{dramati-}}
de \deleted{\unsure{dram}} ``drame'' antique \deleted{\unsure{dans la t}} repris dans la trame de notre vie à
nous, quand \handwritten{elles} ne virent à l'angoisse pure du cauchemar. Nous, acteurs
du drame, \inserted{y} ``marchons'' à fond, certes\handwritten{;} quitte\handwritten{,} au réveil\handwritten{,} \deleted{\unsure{de}} \inserted{à} se sentir tout
idiot et à\deleted{\unsure{x}} se dépêcher de penser à autre chose\handwritten{!} \deleted{\unsure{La}} \handwritten{C}e n'est pas nous, ça
c'est sûr, qui ``parlons'' cette langue étrange, la \deleted{``langue-parabole''}
``langue-images'' ou ``\uline{langue}-\uline{parabole}'', la ``\uline{langue}-\uline{vie}''. Il y\handwritten{\unsure{,}} a un Metteur
en Scène, un plus grand que nous, qui la manie comme \deleted{\unsure{sa}} \inserted{\handwritten{une}} langue maternelle\handwritten{,}
dont nous-mê\handwritten{\unsure{\^{}}}me et la substance même de notre vie \deleted{serions} \inserted{formerions} la chair des
mots. Il monte pêle-mêle drames, farces, idylles ou élégies -- mais même
là où le drame bat son plein, je vois \inserted{\handwritten{pourtant}} que \deleted{\unsure{Celui qui fait jouer fait le}} \inserted{le Narrateur invisible} \deleted{qui fait le récit, il}
garde un sourire en coin\handwritten{\unsure{.|!}} Il sait que la souffrance et la
mort s\handwritten{ont} des choses \inserted{\handwritten{toutes}} simples\handwritten{,} \deleted{et importantes, et} qui n'ont rien de ``dramatique''.
\deleted{\unsure{xx}} Le ``drame'', c'est les vagues que nous \inserted{\handwritten{aimons}} \deleted{\unsure{faisons}} fai\handwritten{re} autour de ces choses\handwritten{,}
histoire de \inserted{\handwritten{nous}} les rendre compliquées...

C'est de \deleted{me rappeller} \inserted{\handwritten{sentir}} ce ``sourire en coin'' du Narrateur, qui fait
que finalement, je \deleted{n'arrive pas \unsure{t}} n'arrive pas à \deleted{\unsure{être entièrement}} \inserted{me sentir} à l'aise
avec \deleted{\unsure{tout un} ``} \handwritten{c}e nom ``psychodrame''\handwritten{,} \deleted{qui ne manque pas} Il a beau avoir de
la gue\deleted{\unsure{l}}\unsure{u|l}le -- je préfère \deleted{\unsure{parler de}} \inserted{\handwritten{le nom serein de}} ``parabole''\handwritten{.} \deleted{aux connotations plus paisibles}
Une parabole \deleted{qui est} ``mise en scène'', c'est une chose entendue\handwritten{,} \deleted{Avec} \inserted{\handwritten{avec un ``scénario'' plus ou moins simple ou plus ou moins imbriqué.}}
\deleted{une} \handwritten{M}ise en scène souvent cocasse, toujours im\unsure{pré|pér}vue et toujours incisive
(sans nul souci, faut il croire, si elle sera comprise...). Elle n'est
limitée par aucune convention\handwritten{,} de style \handwritten{ou} de biènséance ($^{x}$), ni par \handwritten{aucune}
limitation\deleted{s} technique\deleted{s} pour les montages à grand spectacle -- des montages
à faire pâlir les réalisations les plus fantastiques du \handwritten{\unsure{magi}}cien-cinéaste
le plus inspiré\handwritten{!} Si tel ``magicien'' de l'image et \deleted{\illegible{1}} du son parfois nous enchante
ou nous touche, c'est pour avoir su écouter en lui-même, sûrement,
ce ma\handwritten{\unsure{î}}tre-Magicien aux moyens prodigieux, en oeuvre en chacun de nous, et
que si souvent nous dédaignons d'écouter. Il est vrai que ses spectacles
sont gratuits, et qu'il ne \handwritten{f}ait pas de publicité\handwritten{\unsure{.|!}}

\noindent\rule{.23\linewidth}{.3pt}
% slanted
\margin{\handwritten{\unsure{espacer|espace}}}

Ce qui \handwritten{s}ans doute déconcerte le plus dans la langue-paraboles, la
langue-images, c'est la \mbox{l\,i\,b\,e\,r\,t\,é}. On sent que c'est là son âme et
son essence\handwritten{,} \deleted{\unsure{qui}} \handwritten{Elle} a de quoi effrayer plus d'un : une liberté créatrice infinie.
Aucune règle à quoi se raccrocher, ni pour la ``parler'' (mais la
question ne se pose guère -- c'est \deleted{\unsure{quelqu}} un \deleted{\unsure{a}} \handwritten{A}utre en nous qui la parle\handwritten{..)\unsure{,}}
ni pour l'entendre et la saisir -- au vol ! Aucun pas ne semble prescrire
le suivant, ni être prescrit par celui qui l'a précédé -- \deleted{et pourtant,}
\scanpage{19}
% archive sheet number: 18
% slanted
\margin{\unsure{Notes|Note|Nov.} 1982}
$\varepsilon_0,\varepsilon_1\in H$ un groupe OPS $\varepsilon_0,\varepsilon_1$ engendrant
posons
\[\left\{\begin{aligned}
\sigma&=\varepsilon_0\varepsilon_1\varepsilon_0\qquad\text{(\unsure{dissymétrique|dissymétriques})}\\
\rho'_0&=\sigma \varepsilon_0,\quad \rho'_1=\sigma \varepsilon_1\\
\rho_0^{\deleted{\unsure{\prime}}}&=\sigma^{-1}\varepsilon_0\ \deleted{\unsure{\rho_1\varepsilon_0}}\ =\varepsilon_0^{-1}\varepsilon_1^{-1}=(\varepsilon_1\varepsilon_0)^{-1},\quad \rho_1=\deleted{\unsure{\sigma}}\varepsilon_1^{-1}\varepsilon_0^{-1}\deleted{\unsure{\varepsilon}}
\end{aligned}\right.\]
\margin{$\deleted{\unsure{\varepsilon\sigma \varepsilon}\quad \varepsilon_1\varepsilon_0\varepsilon_1^{-1}\stackrel{?}{=}\varepsilon_1}$\\ $\varepsilon_0\varepsilon_1\varepsilon_0\varepsilon_1^{-1}\varepsilon_0^{-1}=\varepsilon_1$\\ On suppose $\varepsilon_0\varepsilon_1\varepsilon_0=\underbrace{\varepsilon_1\varepsilon_0\varepsilon_1}_{\sigma}$}
\[\sigma\ \deleted{\unsure{=|\rho_1}}\ \varepsilon_0\rho_0=\rho_1\varepsilon_0=\rho'_0\varepsilon_0^{-1}=\rho'_1\varepsilon_1^{-1}.\qquad\text{donc }\rho_1=\varepsilon_0(\rho_0)\]
on a alors tjrs
\[\left\{\begin{aligned}
\rho'_1&=\rho_1^{-2}\quad\inserted{\unsure{\varepsilon_1^{-1}}}\\
\sigma(\varepsilon_0)&=\rho'_0(\varepsilon_0)\ \deleted{=\rho_1(\varepsilon_0)}=\deleted{\unsure{\varepsilon_1\varepsilon_0}}\\
&\deleted{\unsure{\rho_0(\varepsilon_0)=}}
\end{aligned}\right.\]
\noindent\rule{\linewidth}{.4pt}

\diagram{$\begin{array}{ccc}(\rho_0,\rho'_0)&& (\rho_1,\rho'_1)\\ &\sigma&\\ \varepsilon_0&&\varepsilon_1\end{array}$}
% slanted
\margin{symétrie conj. par $\sigma$}
\[\text{\textcircled{a}}\quad \sigma(\varepsilon_0)=\varepsilon_1,\quad \sigma(\varepsilon_1)=\varepsilon_0\]
\[\left\{\begin{aligned}\sigma(\rho_0)&=\rho_1,\quad \sigma(\rho_1)=\rho_0,\ \deleted{\unsure{\sigma}}\\\sigma(\rho'_0)&=\rho'_1,\quad \sigma(\rho'_1)=\rho'_0\end{aligned}\right.\]
\margin{$\deleted{\sigma^2(\varepsilon_0)=\varepsilon_0,\ \sigma^2(\varepsilon_1)=\varepsilon_1}$\\ $\sigma^2$ commute à $\deleted{\unsure{\sigma^2(\rho_0}}$ $\varepsilon_0,\varepsilon_1,\rho_0,\rho'_0,\rho_1,\rho'_1$ soit $\sigma^2=\omega$}
% slanted
\margin{\textcircled{$a'$} \uline{conjug.}}
\[\left\{\begin{aligned}
\rho_0(\varepsilon_0)&=\rho'_0(\varepsilon_0)=\varepsilon_1\\
\rho_1(\varepsilon_1)&=\rho'_1(\varepsilon_1)=\varepsilon_0\\
\varepsilon_0(\rho_0)&=\rho_1,\quad \varepsilon_1(\rho_1)=\rho_0\\
\varepsilon_0(\rho'_0)&=\rho'_1,\quad \varepsilon_1(\rho'_1)=\rho'_0
\end{aligned}\right.\]
% slanted
\margin{définition des $\rho_i$ en termes des $\sigma,\varepsilon_i$}
\[\text{\textcircled{b}}\ \left\{\begin{aligned}
\rho_0&=\sigma^{-1}\varepsilon_0\inserted{=\varepsilon_1\sigma^{-1}},\quad \rho_1=\sigma^{-1}\varepsilon_1=\varepsilon_0\sigma^{-1}\\
\rho'_0&=\sigma \varepsilon_0\inserted{=\varepsilon_1\sigma},\quad \rho'_1=\sigma \varepsilon_1=\varepsilon_0\sigma
\end{aligned}\right.\]
% slanted
\margin{introduction de $\omega$\\ éléments d'ordre fini}
\[\text{\textcircled{\unsure{$b'$|b|c}}}\quad \rho_0^3=\rho_1^3=\unsure{\omega|\sigma}^2=\omega,\qquad \omega^2=1\]
\[\left[\ \text{\textcircled{c}}\ \left\{\begin{aligned}\rho_0^{\deleted{\unsure{\prime|1}}6}&=\rho_1^{\unsure{\sigma|6}}=\sigma^4=1\\{\rho'_0}^{3}&={\rho'_1}^{3}=1\end{aligned}\right.\quad \deleted{\unsure{\rho'_0=\rho_0^{-2}}}\ \right]\]
\margin{valables dans $\mathrm{Sl}(2,\mathbb{Z})$ mais non ds $\widetilde{\mathrm{Sl}}(2,\mathbb{Z})$}
% slanted
\margin{relations entre $\rho_i,\rho'_i$}
\[\text{\textcircled{d}}\ \left\{\begin{aligned}
&\deleted{\unsure{\rho'_1=\omega\rho_0^{-2},\quad \rho'_1=\rho_1^{-2}}}\\
\rho'_0&=\omega\rho_0,\quad \rho'_1=\omega\rho_{\deleted{\unsure{2}}1}\\
&\deleted{\unsure{\rho_1=\omega\rho'_0,\quad \rho_1=\omega\rho'_1}}\\
\rho_0^{\unsure{\prime|\vphantom{\prime}}}&=\rho_0^{-2},\quad \rho'_1=\rho_1^{-2}
\end{aligned}\right.\]
\margin{d'où $\rho_0=\omega^{-1}\rho'_0,\ \rho_1=\omega^{-1}\rho'_1,$}
% slanted
\margin{expressions en $\varepsilon_0,\varepsilon_1$}
\[\text{\textcircled{\textcircled{f}}}\ \left\{\begin{aligned}
\sigma&=\varepsilon_0\varepsilon_1\varepsilon_0\ [=\varepsilon_1\varepsilon_0\varepsilon_1]\\
\rho_0&=\varepsilon_0^{-1}\varepsilon_1^{-1}=(\varepsilon_1\varepsilon_0)^{-1}\unsure{-},\quad \rho_1=\varepsilon_1^{-1}\varepsilon_0^{-1}=(\varepsilon_0\varepsilon_1)^{-1}\\
\rho_0^{-1}&=\varepsilon_1\varepsilon_0,\qquad \rho_1^{-1}=\varepsilon_0\varepsilon_1
\end{aligned}\right.\]
% slanted
\margin{expressions en $\sigma,\varepsilon_0$}
\[\text{\textcircled{\textcircled{g}}}\ \left\{\begin{aligned}
\varepsilon_1&=\sigma(\varepsilon_0)=\sigma \varepsilon_0\sigma^{-1}\\
\rho_0&=\sigma^{-1}\varepsilon_0,\quad \rho'_0=\sigma \varepsilon_0\\
\rho_1&=\varepsilon_0^{\unsure{1|\vphantom{1}}}\sigma^{-1},\quad \rho'_1=\deleted{\unsure{\sigma}}\varepsilon_0\sigma
\end{aligned}\right.\]
% slanted
\margin{expressions en $(\sigma,\rho_0)$}
\[\text{\textcircled{\textcircled{h}}}\ \left\{\begin{aligned}
\varepsilon_0&=\sigma\rho_0\quad\deleted{\unsure{\varepsilon_1=}}\\
\varepsilon_1&=\rho_0\sigma\\
\rho'_0&=\rho_0^{-2}\\
\rho_1&=\sigma(\rho_0)=\sigma\rho_0\sigma^{-1}\\
\rho'_1&=\sigma(\rho_0^{-2})=\sigma\rho_0^{-2}\sigma^{-1}
\end{aligned}\right.\]
% slanted
\margin{expressions en $(\sigma,\rho'_0)$}
\[\text{\textcircled{\textcircled{i}}}\ \left\{\begin{aligned}
\varepsilon_0&=\sigma^{-1}\rho'_0\\
\varepsilon_1&=\rho'_0\sigma^{-1}\\
\rho_0&=\sigma^2\rho'_0=\rho'_0\sigma^2\\
\rho_1&=\sigma^{-1}\rho'_0\sigma^{-1}\ \deleted{\unsure{\omega\sigma\rho'_0\sigma}}\\
\rho'_1&=\sigma(\rho'_0)=\sigma\rho'_0\sigma^{-1}
\end{aligned}\right.\]
% slanted
\margin{expressions en $(\varepsilon_0,\rho_0)$}
\[\text{\textcircled{j}}\ \left\{\begin{aligned}
\varepsilon_1&=\rho_0(\varepsilon_0)=\rho_0\varepsilon_0\rho_0^{-1}\\
\sigma&=\unsure{\rho_0}\varepsilon_0\rho_0^{-1}\\
\rho'_0&=\rho_0^{-2}\\
\rho_1&=\varepsilon_0(\rho_0)=\varepsilon_0\rho_0\varepsilon_0^{-1}\\
\rho'_1&=\varepsilon_0(\rho_0^{-2})=\varepsilon_0\rho_0^{-2}\varepsilon_0^{-1}
\end{aligned}\right.\]
\scanpage{20}
% archive sheet number: 19
% slanted
\margin{expressions en $(\varepsilon_0,\rho'_0)$}
\[k)\ \left\{\begin{aligned}
\varepsilon_1&=\rho'_0(\varepsilon_0)=\rho'_0\varepsilon_0{\rho'_0}^{-1}\\
\sigma&=\rho'_0\varepsilon_0^{-1}\\
\rho_0&=\varepsilon_0{\rho'_0}^{-1}\varepsilon_0\\
\rho_1&=\varepsilon_0^2{\rho'_0}^{-1}\\
\rho'_1&=\varepsilon_0^{\deleted{\unsure{-1}}}(\rho'_0)=\varepsilon_0\rho'_0\varepsilon_0^{-1}
\end{aligned}\right.\]
% slanted
\margin{expressions en $(\varepsilon_0,\rho_1)$}
\[\text{\textcircled{l})}\ \left\{\begin{aligned}
\varepsilon_1&=\rho_1^{-1}(\varepsilon_0)=\rho_1^{-1}\varepsilon_0\rho_1\\
\sigma&=\rho_1^{-1}\varepsilon_0\\
\rho_0&=\varepsilon_0^{-1}(\rho_1)=\varepsilon_0^{-1}\rho_1\varepsilon_0\\
\rho'_1&=\rho_1^{-2}\\
\rho'_0&=\varepsilon_0^{-1}(\rho_1^{-2})=\varepsilon_0^{-1}\rho_1^{\unsure{-1|-7|-2}}\varepsilon_0
\end{aligned}\right.\]
% slanted
\margin{expressions en $(\varepsilon_0,\rho'_1)$}
\[\text{\textcircled{\unsure{m}}}\ \left\{\begin{aligned}
\deleted{\unsure{\sigma=\varepsilon_0^{-1}\rho'_1}}\quad \varepsilon_1&={\rho'_1}^{-1}(\varepsilon_0)={\rho'_1}^{-1}\varepsilon_0\rho'_1\\
\sigma&=\varepsilon_0^{-1}\rho'_1\\
\rho_1&=\varepsilon_0{\rho'_1}^{-1}\varepsilon_0\\
\rho'_0&=\varepsilon_0^{-1}(\rho'_1)=\varepsilon_0^{-1}\rho'_1\varepsilon_0\\
\rho_0&={\rho'_1}^{-1}\varepsilon_0^2
\end{aligned}\right.\]
% slanted
\margin{19.6.83}

\uline{NB} parmi les sept éléments $\rho_i,\rho'_i,\varepsilon_i,\sigma$, \deleted{deux \unsure{q}} \inserted{une}
paire quelconque qui n'est pas formée uniquement
d'éléments parmi les 4 $\rho$ [ce qui fait $\binom{7}{2}-\binom{4}{2}$
$=21-6=15$ paires) engendrent $H$. D'autre part,
les quatre éléments $\unsure{\rho_i|\rho}$ $(\rho_0,\rho_1,\rho'_0,\rho'_1)$ n'engendrent
pas $H=\langle \varepsilon_0,\varepsilon_1\mid \varepsilon_0\varepsilon_1\varepsilon_0=\varepsilon_1\varepsilon_0\varepsilon_1\rangle$, \unsure{par} $\unsure{=|\simeq}$ mod commutateurs\unsure{...}.

Il faudrait aussi regarder l'extension canonique
de $\mathbb{Z}/2\mathbb{Z}$ par
\[\begin{aligned}
\unsure{S\mathcal{T}|S\mathscr{T}}_{1,1}^+&=\langle \varepsilon_0,\varepsilon_1\mid \varepsilon_0\varepsilon_1\varepsilon_0=\varepsilon_1\varepsilon_0\varepsilon_1\rangle\\
&\simeq\langle\rho,\sigma\mid \rho^{-3}=\sigma^2\rangle\\
&\simeq\langle\rho,\varepsilon_0\mid \rho_{\unsure{0|\vphantom{0}}}\varepsilon_0\rho_{\unsure{0|\vphantom{0}}}^{-1}\varepsilon_0\rho=1\rangle
\end{aligned}\]
donnant naissance (par \unsure{division} par la relation $\omega=1$)
à $\mathrm{Gl}(2,\mathbb{Z})$ ext. de $\mathbb{Z}/2\mathbb{Z}$ par $\mathrm{Sl}(2,\mathbb{Z})$.
\scanpage{21}
% archive sheet number: 20
\unsure{non} l'égalité -- celle-ci sert à \unsure{sommer} (72).

La situation devient encore plus jolie si on suppose
$\unsure{\zeta|z}=1$, puisque alors la donnée supplémentaire
revient à la donnée
(\unsure{69|09|0g|eq})\qquad un sous-groupe \inserted{(commutatif)} $T_H$ de $H$, normalisant $U_0$,
tel que $T_H\cap U_0=\{1\}$, et que $\sigma(t)=t^{-1}$ pour
$t\in T_H$.
% A long diagonal runs across the square and lower working, ending at the last braid word before the two final tests.
\deleted{\diagram{\begin{tikzcd}[ampersand replacement=\&]
\widetilde Z\arrow[r,"\unsure{k|K|\varepsilon}"]\arrow[d,"u_0"'," \unsure{2|\wr}"] \& Z\arrow[d,"u"," \unsure{2|\wr}"']\\
T_0\arrow[r,"c_0"] \& T_H
\end{tikzcd}}
\diagram{$\begin{array}{cc}\lambda^{-1}&\lambda^{\unsure{\omega|m}}\\1&\lambda^2\end{array}$}
\[\widetilde Z\]

$\varepsilon_0,\varepsilon_1,\sigma,\rho_0,\rho_1,\rho'_0,\rho'_1,\omega$
\margin{$\rho_0^3=\rho_1^3=\sigma^2=\omega$\\ $\omega^2=1$, $\omega$ commute à tout}
\[\left\{\begin{aligned}
\rho'_0&=\sigma \varepsilon_0 &\rho_0&=\inserted{\sigma^{-1}\varepsilon_0=}(\varepsilon_1\varepsilon_0)^{-1} &\rho'_0&=\inserted{\deleted{\unsure{\omega\rho\varepsilon_0}}}\rho_0^{-2} &{\rho'_0}^3&=1 &\rho_0^6&=1\\
\rho'_1&=\sigma \varepsilon_1 &\rho_1&=\inserted{\sigma^{-1}\varepsilon_1=}(\varepsilon_0\varepsilon_1)^{-1} &\rho'_1&=\rho_1^{-2} &{\rho'_1}^3&=1 &\rho_1^6&=1
\end{aligned}\right.\]
\margin{$\rho'_0=\omega\rho_0,\ \rho_0=\omega\rho'_0$\\ $\rho'_1=\omega\rho_1,\ \rho_1=\omega$}
\[\left\{\begin{aligned}
\varepsilon_1&=\sigma(\varepsilon_0)=\rho_{\unsure{1|0}}(\varepsilon_0)=\rho'_0(\varepsilon_0)=\varepsilon_1(\varepsilon_0)\\
\varepsilon_0&=\sigma(\varepsilon_1)=\rho_0(\varepsilon_1)=\rho'_1(\varepsilon_1)=\varepsilon_0(\varepsilon_1)\\
\sigma&=\varepsilon_1\varepsilon_0\varepsilon_1=\varepsilon_0\varepsilon_1\varepsilon_0
\end{aligned}\right.\tag*{$(\ast)$}\]

Partons de $\varepsilon_0,\varepsilon_1$ [satisfaisant $(1)\ \boxed{\varepsilon_1\varepsilon_0\varepsilon_1=\varepsilon_0\varepsilon_1\varepsilon_0}$]
posons $\sigma=\varepsilon_0\varepsilon_1\varepsilon_0[=\varepsilon_1\varepsilon_0\varepsilon_1],\rho'_0=\sigma \varepsilon_0,\rho'_1=\sigma \varepsilon_1,\rho_0^{\unsure{?|-1}}=\varepsilon_1\varepsilon_0,\rho_1^{\unsure{1|-1}}=\varepsilon_0\varepsilon_1$
\deleted{je dis qu'on a bien} \inserted{examinons les formules $(\ast)$} \deleted{\unsure{alors}} qui restent à vérifier. On a
\[\underbrace{\sigma(\varepsilon_0)=\rho'_0(\varepsilon_0)}=\varepsilon_0\varepsilon_1(\varepsilon_0)\deleted{\unsure{\varepsilon_0\varepsilon_0^{-1}\varepsilon_1^{-1}\varepsilon_1^{-1}}}=\varepsilon_0\varepsilon_1\varepsilon_0\varepsilon_1^{-1}\varepsilon_0^{-1}\stackrel{?}{=}\varepsilon_1\]
c'est égal à $\varepsilon_0$ ssi \deleted{$\varepsilon_1\varepsilon_0=\varepsilon_0\varepsilon_1$} on a (1)
\[\sigma(\varepsilon_1)=\rho'_1(\varepsilon_1)=\varepsilon_0\underline{\varepsilon_1\varepsilon_0\varepsilon_1}\varepsilon_0^{-1}\varepsilon_1^{-1}\varepsilon_0^{-1}\stackrel{?}{=}\varepsilon_0\]
égal à $\varepsilon_0$ ssi $\varepsilon_1\varepsilon_0\varepsilon_1=\varepsilon_0\varepsilon_1\varepsilon_0$ i.e (1)}
\[\begin{aligned}
\rho_0^2&=\varepsilon_1\varepsilon_0\varepsilon_1\varepsilon_0\stackrel{?}{=}\rho'_0=\varepsilon_0\varepsilon_1\varepsilon_0\mathbin{\cdot}\varepsilon_0\quad\text{ssi on a (1)}\\
\rho_1^2&=\varepsilon_0\varepsilon_1\varepsilon_0\varepsilon_1\stackrel{?}{=}\rho'_1=\varepsilon_0\varepsilon_1\varepsilon_0\varepsilon_1\quad\text{tjrs vrai}
\end{aligned}\]
\scanpage{22}
% archive sheet number: 21
\[\underbrace{\begin{pmatrix}0&\beta\\-\beta^{-1}&0\end{pmatrix}\begin{pmatrix}1&\beta\\0&1\end{pmatrix}}_{\begin{pmatrix}0&\beta\\-\beta^{-1}&-1\end{pmatrix}}\begin{pmatrix}0&\beta\\-\beta^{-1}&0\end{pmatrix}=\begin{pmatrix}-1&0\\\beta^{-1}&-1\end{pmatrix}\]
\margin{$\unsure{v|V|\nu}_0(\beta)^{-1}=$}
\[g=\begin{pmatrix}\unsure{a|\mu}&b\\c&d\end{pmatrix}\qquad g^{-1}=\begin{pmatrix}d&-b\\-c&a\end{pmatrix}\]
\margin{$1-a^2c^2\beta^2$}
\[\begin{aligned}
&\underbrace{g\begin{pmatrix}1&\beta\\0&1\end{pmatrix}}\deleted{\unsure{=}}g^{-1}\\
&\begin{pmatrix}a&a\beta+b\\c&c\beta+d\end{pmatrix}\begin{pmatrix}d&-b\\-c&a\end{pmatrix}=\begin{pmatrix}1-ac\beta&a^2\beta\\-c^2\beta&1+ac\beta\end{pmatrix}
\end{aligned}\]
\[\deleted{e_1(\beta)}\qquad \deleted{e_1(\gamma)e_0(-\gamma^{-1})e_1(\gamma)=}\]
\margin{$ad-bc$}
\[\underbrace{v_0(\beta)e_0(\beta)}\underbrace{v_0(\beta)^{-1}e_0(\beta)}\underbrace{v_0(\beta)e_0(\beta)}v_0(\beta^{\deleted{\unsure{-1}}})^{-1}=v_0(\beta)\]
\[\omega\big(v_0(\beta)e_0(\beta)\big)^3=\]
% margin note written sideways
\margin{$\rho(\beta)=\begin{pmatrix}0&\beta\\-\beta^{-1}&0\end{pmatrix}\begin{pmatrix}1&\beta\\0&1\end{pmatrix}=\begin{pmatrix}0&\beta\\-\beta^{-1}&-1\end{pmatrix}=u(\unsure{\beta})\bigl($}
% margin note written sideways
\margin{$\begin{pmatrix}0&1\\-1&-1\end{pmatrix}\begin{pmatrix}0&1\\-1&-1\end{pmatrix}=\begin{pmatrix}-1&-1\\1&0\end{pmatrix}$}
% margin note written sideways
\margin{$ {\rho'}^3=\begin{pmatrix}-1&-1\\1&0\end{pmatrix}\begin{pmatrix}0&1\\-1&-1\end{pmatrix}=\left(\begin{matrix}1&0\\0&2\end{matrix}\right]$}
\[\underbrace{\begin{pmatrix}1&0\\\gamma&1\end{pmatrix}\begin{pmatrix}1&\beta\\0&1\end{pmatrix}}_{\begin{pmatrix}1&\beta\\\gamma&\beta\gamma+1\end{pmatrix}}\begin{pmatrix}1&0\\\gamma&1\end{pmatrix}=\begin{pmatrix}1+\beta\gamma&\beta\\\underset{=\gamma(2+\beta\gamma)}{2\gamma+\beta\gamma^2}&1+\beta\gamma\end{pmatrix}\]
\[1+\beta^2\gamma^2+2\beta\gamma-2\beta\gamma-\beta^2\gamma^2\]
\scanpage{23}
% archive sheet number: 22
\[\{s,\qquad \unsure{v|V|\nu}_0(\beta)\deleted{\unsure{\bigl(}}\big(e_0(\beta)\big)=e_1(-\beta^{-1})\]
\[\begin{pmatrix}0&\beta\\-\beta^{-1}&0\end{pmatrix}\begin{pmatrix}\unsure{1|0}&\beta\\0&1\end{pmatrix}\bigl(\]
% The following draft is turned through 180 degrees and crossed by long diagonal strokes.
\deleted{Remarque considérons un multigraphe \deleted{$\unsure{(S,(A_i)_{i\in\Pi },Z\subset G_S)|(S,(A_\pi)_{\pi\in\Pi},Z\subset\mathfrak S_S)}$}
\deleted{$\unsure{M=(\Gamma|\mathcal M=(}$} $M=(\underset{\substack{\Vert\\(S,\check S,R\subset S\times\check S)}}{\Gamma},\underset{\substack{\text{\unsure{ss-groupe|ss-groupes}}\\\text{comm.}}}{Z\subset\mathrm{Aut}(\Gamma)},\underset{\substack{\Vert\\ Z\text{-torseurs}}}{\Pi}\subset\Pi _Z\subset\mathrm{Bij}(S,\check S))$
avec op. libres de $Z$ sur $S$ et $\check S$,
il lui est donc associé

a) Un ens. $T$ $(\simeq S^{\unsure{\leftrightarrows|/Z}}\simeq\check S^{\unsure{\leftrightarrows|/Z}})$

b) \deleted{Deux familles de} Un groupe $Z$ commutatif

c) Deux \deleted{familles \unsure{de}} \inserted{familles de} $Z$-torseurs, \deleted{sur $T$}, indexées
par $T$ \deleted{et en \unsure{deux}} (savoir les fibres de $t$ dans $S$ et dans $\check S$ respectivement, i.e. les \unsure{voyeurs|voyous|origines} et \unsure{voyeuses|voyouses|extrémités})
\margin{$t$\\ où}
\deleted{\unsure{d)}} \unsure{en tant}\qquad $Z$
% margin note written sideways
\margin{$\underline{(s,\lambda t)}$}
% margin note written sideways
\margin{$\lambda s_0$}
\[(\unsure{s_0},\check s_0),\quad(s_0,\lambda\check s_0)\in R\]
\[\exists g\in G_0\qquad g(s_0)=s_0,\quad g(\check s_0)=\lambda\check s_0\]
\[g^{-1}(s_0)=s_0)\quad\circ\quad \unsure{g'|g^{-1}}\lambda g^{-1}(s_0)=\]
\margin{$\lambda\{s,t\}$ $\in A_{\unsure{\pi|i}}$\\ $\{\lambda s,t\}$ $\in A_{\unsure{\pi|i}}$\\ $\{s,\lambda t\}$ $\in A_{\unsure{\pi|i}}$\\ $t,\lambda^{-1}s\lambda$\\ \deleted{\unsure{(s)}}}
\[g s_0\qquad \lambda s_0=s,\quad \lambda^{-1}t_0=t_0\]
\[U g^{\unsure{\prime}}\qquad t.u_0=\deleted{\unsure{g}}t'u'_0\deleted{\unsure{g}}=\sigma({t'}^{-1}u'_0)\]
\[\deleted{t=u_0u_0^{-1}}\qquad T\in\]
\[\deleted{\unsure{u\in U_1|u\in U_0}}\qquad \deleted{u\in U_0}\qquad T\cap U_0\]
\[us_0=\lambda s_0\qquad ut_0=t_0\]
\[u\in T\mathbin{\cdot}U_0\quad u\in U_1\qquad \boxed{TU_0\cap U_1=\unsure{\zeta|z|\{1\}}}\]
\[TU_0\cap U_1=\zeta\quad(\unsure{72|12})\quad\deleted{\Rightarrow T\unsure{u_0t_0|\text{ normalise }}\cap U_1}\qquad \Rightarrow B_0\cap B_1=T\]
\[B_0\cap B_1\qquad T\cap U_0\underset{?}{\subset}U_1\]}
\scanpage{24}
% archive sheet number: 23
$M=L\oplus L'$\qquad \deleted{\unsure{A et}} $L,L'$ $\Lambda$-modules inversibles
\[\underbrace{\mathrm{Sl}(M)}_{\unsure{H}}=\left\{\begin{pmatrix}a&b\\c&d\end{pmatrix}\ \middle|\ \begin{aligned}&a,d\in\Lambda\\&c\in L^{-1}L',b\in{L'}^{-1}L\end{aligned}\ \deleted{\unsure{\big\}}},\ ad-bc=1\right\}\]
\[U_0=\left\{\begin{pmatrix}1&\alpha\\0&1\end{pmatrix}\ \middle|\ \alpha\in\overset{L_0}{\overbrace{\underbrace{{L'}^{-1}L}_{\deleted{\unsure{L_1}}}}}\right\}\]
\[B_0=\left\{\begin{pmatrix}\lambda&\alpha\\0&\lambda^{-1}\end{pmatrix}\ \middle|\ \lambda\in\Lambda^*,\alpha\in{L'}^{-1}L\right\}\]
\[U_1=\left\{\begin{pmatrix}1&0\\\beta&1\end{pmatrix}\ \middle|\ \beta\in\underbrace{L^{-1}L'}_{L_1}\right\}\]
\[B_1=\left\{\begin{pmatrix}\lambda&0\\\beta&\lambda^{-1}\end{pmatrix}\ \middle|\ \lambda\in\Lambda^{\unsure{\times|\ast}},\deleted{\unsure{\big\}}}\beta\in L^{-1}L'\right\}\]
\[T=\left\{\begin{pmatrix}\lambda&0\\0&\lambda^{-1}\end{pmatrix}\ \middle|\ \lambda\in\Lambda^*\right\}\]
\[\begin{pmatrix}a&b\\c&d\end{pmatrix}^{-1}=\begin{pmatrix}d&-b\\-c&a\end{pmatrix}\]
\[\sigma=\begin{pmatrix}0&-1\\1&0\end{pmatrix}\quad\text{si }L=L'=\Lambda\]
\[\begin{aligned}
\deleted{\sigma(c)}\ \deleted{\unsure{N(c)=\begin{pmatrix}0&-c\\1/c&0\end{pmatrix}}}\quad N(c)&=\begin{pmatrix}0&-c^{-1}\\c&0\end{pmatrix}\\&=\sigma\cdot u(c)
\end{aligned}\]
si $g=\begin{pmatrix}a&b\\c&d\end{pmatrix}$, on a
\[\begin{aligned}
g\begin{pmatrix}\lambda&0\\0&\lambda^{-1}\end{pmatrix}&=\begin{pmatrix}a\lambda&b\lambda^{-1}\\c\lambda&d\lambda^{-1}\end{pmatrix}\\
\begin{pmatrix}\lambda&0\\0&\lambda^{-1}\end{pmatrix}g&=\begin{pmatrix}\lambda a&\lambda b\\\lambda^{-1}c&\lambda^{-1}d\end{pmatrix}\\
g\begin{pmatrix}\lambda&0\\0&\lambda^{-1}\end{pmatrix}g^{-1}&=\begin{pmatrix}a&b\\c&d\end{pmatrix}\begin{pmatrix}\lambda d&-\lambda b\\-\lambda^{-1}c&\lambda^{-1}a\end{pmatrix}\\
&=\begin{pmatrix}\lambda ad-\lambda^{-1}bc&\deleted{\unsure{a\lambda b}}ab(-\lambda+\lambda^{-1})\\cd(\lambda-\lambda^{-1})&-\lambda bc+\lambda^{-1}ad\end{pmatrix}
\end{aligned}\]
\[\begin{aligned}
g&=\begin{pmatrix}a&b\\c&d\end{pmatrix}\\u(\lambda)&=\begin{pmatrix}\lambda&0\\0&\lambda^{-1}\end{pmatrix}\\e(\alpha)&=\begin{pmatrix}1&\alpha\\0&1\end{pmatrix}\\e'(\beta)&=\begin{pmatrix}1&0\\\beta&1\end{pmatrix}
\end{aligned}\]
\[\begin{aligned}
N(c)u(\lambda)&=N(c\lambda)\\u(\lambda)N(c)&=N(\lambda^{-1}c)\\u(\lambda)N(c)u(\lambda')&=N(\lambda^{-1}\lambda'c)
\end{aligned}\]
\[\begin{aligned}
N&=\left\{\begin{pmatrix}a&b\\c&d\end{pmatrix}\in\mathrm{Sl}(M)\ \middle|\ ab\deleted{\unsure{(\lambda-\lambda^{-1})}}=cd=0\right\}\\
\deleted{\unsure{N^+}}\quad N^+&=T=\left\{\begin{pmatrix}a&b\\c&d\end{pmatrix}\in\mathrm{Sl}(M)\ \middle|\ b=c=0\right\}\\
N^-\,[=\sigma T]&=\left\{\begin{pmatrix}a&b\\c&d\end{pmatrix}\in\mathrm{Sl}(M)\ \middle|\ a=d=0\right\}
\end{aligned}\]
\[\begin{aligned}u(\lambda)\sigma\deleted{\unsure{u_\lambda}}&=\begin{pmatrix}0&-\lambda\\\lambda^{-1}&0\end{pmatrix}\\\sigma u(\lambda)&=\begin{pmatrix}0&-\lambda^{-1}\\\lambda&0\end{pmatrix}\end{aligned}\]
\[\boxed{\begin{aligned}
e(\alpha)u(\lambda)e'(\beta)&=\begin{pmatrix}1&\alpha\\0&1\end{pmatrix}\underbrace{\begin{pmatrix}\lambda&0\\0&\lambda^{-1}\end{pmatrix}\begin{pmatrix}1&0\\\beta&1\end{pmatrix}}_{\begin{pmatrix}\lambda&0\\\lambda^{-1}\beta&\lambda^{-1}\end{pmatrix}}=\begin{pmatrix}\lambda+\lambda^{-1}\alpha\beta&\lambda^{-1}\alpha\\\lambda^{-1}\beta&\lambda^{-1}\end{pmatrix}=\begin{pmatrix}a&b\\c&d\end{pmatrix}\\
&\deleted{\unsure{d=\lambda^{-1}}}\qquad \lambda=d^{-1}\qquad \alpha=\lambda\unsure{\beta|b}=b/d\qquad \beta=\lambda c=c/d
\end{aligned}}\]
donc $U_0\times T\times U_1\to{H}$ injectif
\[U_0TU_1=\left\{\begin{pmatrix}a&b\\c&d\end{pmatrix}\in{H}\ \middle|\ d\in\Lambda^*\right\}\]
\[\begin{aligned}
[e(\alpha),e'(\beta)]&=\underbrace{e(\alpha)e'(\beta)}e(-\alpha)e'(-\beta)=\begin{pmatrix}(1+\alpha\beta)^2-\alpha\beta&-\alpha^2\beta\\\alpha\beta^2&1-\alpha\beta\end{pmatrix}=\begin{pmatrix}1+\alpha\beta+\alpha^2\beta^2&-\alpha^2\beta\\\alpha\beta^2&1-\alpha\beta\end{pmatrix}\\
&\begin{pmatrix}1+\alpha\beta&\alpha\\\beta&1\end{pmatrix}\begin{pmatrix}1+\alpha\beta&-\alpha\\-\beta&1\end{pmatrix}
\end{aligned}\]
si $\alpha\beta=1$, on trouve $\deleted{\unsure{=}}[e(\alpha),e'(\beta)]=\begin{pmatrix}3&-\alpha\\\alpha^{-1}&0\end{pmatrix}$ \deleted{\unsure{ça ne p}}
donc si $\exists \lambda$ tel que $[e(\alpha),e'(\beta)]\in N^-$, alors $\alpha\beta=1$, $3=0$
\inserted{i.e. on est en car 3}
\scanpage{25}
% archive sheet number: 24
\deleted{Le calcul (64') donne, si $\lambda\in Z,\tilde \lambda\in\widetilde Z$, avec $\lambda=\tilde \lambda^{\unsure{\bullet|\cdot}}$
\margin{$\Vert$\\ $\delta_{\widetilde Z}(\unsure{\lambda|\tilde \lambda})$}
\deleted{$\unsure{(70)}\left\{\begin{aligned}&\text{Si }\lambda=\deleted{\unsure{\lambda}}\tilde \lambda^\bullet,\ \text{on a}\\&u_{\unsure{\lambda}}^{-1}\unsure{v|\nu}_{\tilde \lambda}^{2}(s_0,t_0\unsure{\,;}\end{aligned}\right\}$}
\[(u_\lambda^{-1}v_{\tilde \lambda}^2)(s_0,t_0)=(\lambda s_0,\lambda t_0)\]
i.e.
\[\lambda^{-1}u_\lambda^{-1}v_{\tilde \lambda}^{2}(s_0,t_0)=(s_0,t_0),\]
ce qui implique que l'on a
\[\lambda^{-1}u_\lambda^{-1}v_{\tilde \lambda}^2\in T'\cap G_0=\unsure{\zeta|\mathfrak z}_{T'}\]
mais on a aussi
\[\widetilde \delta(\lambda^{-1}u_\lambda^{-1}v_{\tilde \lambda}^2)=\widetilde \delta(\lambda^{-1})\underset{1}{\underline{\widetilde \delta(u_\lambda^{-1})}}\underset{\deleted{\unsure{\lambda}}\ \tilde \lambda^2}{\underline{\widetilde \delta(v_{\tilde \lambda}^{\deleted{\unsure{2}}})^2}}\]
je voudrais montrer que c'est $=1$, i.e. qu'on a
\[\widetilde \delta(\lambda)=\tilde \lambda^2\quad\cdot\quad\forall\tilde \lambda\in\widetilde Z,\quad \lambda=\tilde \lambda^\bullet\in Z\tag*{(70)}\]
i.e. que l'on \deleted{\unsure{a que si}} a commutativité dans
\diagram{\begin{tikzcd}[ampersand replacement=\&,column sep=large]
\widetilde Z\arrow[r,"{\delta_{\unsure{\widetilde Z|\overline Z}}:\tilde \lambda\mapsto\tilde \lambda^\bullet}"]\arrow[rr,bend right=35,"{\tilde \lambda\mapsto\tilde \lambda^2}"'] \& Z\arrow[r,"{\widetilde \delta\mid Z}"] \& \widetilde Z
\end{tikzcd}}}
\scanpage{26}
% archive sheet number: 25
\[e'(\beta)e(\alpha)=\begin{pmatrix}1&0\\\beta&1\end{pmatrix}\begin{pmatrix}1&\alpha\\0&1\end{pmatrix}=\begin{pmatrix}1&\alpha\deleted{\unsure{a}}\\\beta&1+\alpha\beta\end{pmatrix}\]
Si $1+\alpha\beta\in\Lambda^*$, on a
\[\boxed{\boxed{e'(\beta)e(\alpha)=e\left(\frac{\alpha}{1+\alpha\beta}\right)u\left(\frac{1}{1+\alpha\beta}\right)e'\left(\frac{\beta}{1+\alpha\beta}\right)}}\]
Si $1+\alpha\beta=0$ i.e $\alpha\beta=-1$ on a
\[e'(\beta)e(\alpha)=\begin{pmatrix}1&\alpha\\\deleted{\unsure{-\alpha}}\beta&0\end{pmatrix}\underset{?}{\in}[\sigma B_1=]N^-U_1\]
% margin note written sideways
\margin{$\deleted{\unsure{-\alpha\beta-\alpha^2\beta^2-\alpha^3\beta^3}}$\\ $\deleted{1-\alpha^3\beta^3+\alpha^3\beta^3}$\\ $\deleted{e'(\beta)e(\alpha)}$}
\[\deleted{\underbrace{\sigma}_{\begin{pmatrix}0&-1\\1&0\end{pmatrix}}\underbrace{\begin{pmatrix}\unsure{0|a'}&a'\\\unsure{-b'}+\lambda^{-1}&\lambda^{-1}\end{pmatrix}\begin{pmatrix}1&0\\\beta'&1\end{pmatrix}}_{\begin{pmatrix}\lambda&0\\\lambda^{-1}\beta'&\lambda^{-1}\end{pmatrix}}=\begin{pmatrix}-\lambda^{-1}\beta'&-\lambda^{-1}\\\lambda&0\end{pmatrix}=\begin{pmatrix}a&b\\c&d\end{pmatrix}}\]
\[\begin{aligned}
\deleted{\unsure{N^-=\{}}\quad N^-B_1&=\left\{\begin{pmatrix}a&b\\c&d\end{pmatrix}\in H\mid d=0\right\}\\
&=\left\{\begin{pmatrix}\unsure{a'|a}&b\\c&0\end{pmatrix}\ \middle|\ \begin{aligned}a\in\Lambda,\quad b\in{L'}^{-1}L,\ c\in L^{-1}L'\\bc=\deleted{\unsure{1}}-1\end{aligned}\right\}
\end{aligned}\]
\[\underbrace{\begin{pmatrix}0&b'\\-{b'}^{-1}&0\end{pmatrix}}_{\in N^-}\underbrace{\begin{pmatrix}1&\overset{\deleted{\unsure{a|=}}}{0}\\\beta'\,\unsure{t?|\alpha}&1\end{pmatrix}}_{\in U_1}=\begin{pmatrix}b'\beta'&b'\\-{b'}^{-1}&0\end{pmatrix}\]
alors
\[\begin{pmatrix}a&b\\c&0\end{pmatrix}=\underbrace{\begin{pmatrix}0&b\\c&0\end{pmatrix}}_{\in N^-}\deleted{\unsure{\frac{a}{cb}}}\underbrace{\begin{pmatrix}1&0\\\underbrace{ab^{-1}}&1\end{pmatrix}}_{\in U_1}\]
si $bc=-1$
\[\deleted{e'(\beta)e(-\beta^{-1})=}\]
Si $\alpha\beta=-1$
\[e'(\beta)e(\alpha)=\underbrace{\begin{pmatrix}0&\alpha\\\beta&0\end{pmatrix}}_{\in N^-}\begin{pmatrix}1&0\\\alpha^{-1}=-\beta&1\end{pmatrix}\qquad\text{i.e.}\]
\[e'(\beta)e(-\beta^{-1})=\begin{pmatrix}0&-\beta^{-1}\\\beta&0\end{pmatrix}\underbrace{\begin{pmatrix}1&0\\-\beta&1\end{pmatrix}}_{e'(\unsure{-\beta|\beta})}\qquad\text{i.e.}\]
\[e'(\beta)e(-\beta^{-1})=N(\beta)e'(-\beta)\]
\[\boxed{\boxed{e'(\beta)e(-\beta^{-1})e'(\beta)=N(\beta)}}\qquad \beta\in(L^{-1}L')^*=U_1^{*}\]
\[\deleted{e'(-\alpha^{-1})e(\alpha)e'(-\alpha^{-1})=N(-\alpha^{-1})}\]
% margin note written sideways
\margin{$1\quad \lambda+\lambda^{-1}\alpha\beta=d^{-1}+\deleted{\unsure{\lambda}}d\frac{bc}{d^2}=d^{-1}+\frac{bc}{d}=a$\\ i.e. $\deleted{\unsure{1+bc=ad}}$\quad OK}
% margin note written sideways
\margin{$\deleted{-\lambda^2abcd-\lambda^{-2}abcd+a^2d^2+b^2c^2+abcd(\lambda^2+\lambda^{\unsure{-2|2}}\deleted{\unsure{+}}2)}$\\ $=\deleted{abcd(\ a^2d^2+b^2c^2-2adbc=(ad-bc)^2=1}$}
\scanpage{27}
% archive sheet number: 26
\mspage{114}
% The following draft is upside down in the scan and cancelled by a long diagonal stroke.
 \deleted{On considère alors le groupe
\[
 \widetilde Z=G/G_0^+\qquad\text{avec}\qquad G\xrightarrow{\tilde\delta}\widetilde Z \tag*{(9)}
\]
qui est un rev. de degré 2 de $Z$
\[
 1\longrightarrow\{\pm1\}\xrightarrow[\unsure{\prime|,}]{}  \widetilde Z\longrightarrow Z\longrightarrow 1,\tag*{(10)}
\]
et le diagramme d'inclusions de groupes
\diagram{\begin{tikzcd}[ampersand replacement=\&,column sep=small,row sep=small]
 {} \& U\arrow[r,hook,no head,"Z","{\deleted{\unsure{\mathrm{dist.}}}}"']\arrow[d,hook,no head,"\mu_2"'] \& B''\arrow[r,hook]\arrow[d,hook,no head,"\mu_2"'] \& G_0^+\arrow[d,hook,no head,"\mu_2"'] \& {}\\
\widetilde Z\arrow[u]\arrow[d] \& \widetilde U\arrow[r,hook,no head,"Z","{\deleted{\unsure{\mathrm{dist.}}}}"']\arrow[d,hook,"Z"'] \& \widetilde B''\arrow[r,hook]\arrow[d,hook,"Z"'] \& G_0\arrow[d,hook,"Z"'] \& \widetilde Z\arrow[u]\arrow[d]\\
{} \& B'\arrow[r,hook,no head,"Z","{\deleted{\unsure{\mathrm{dist.}}}}"'] \& B\arrow[r,hook,no head] \& G \& {}
\end{tikzcd}\quad (11)}
 \margin{On a écrit au dessus de toute flèche d'inclusion d'un sous-groupe dist. dans un groupe \unsure{ambiant}, le groupe quotient correspondant (à isom. can. près)}
 \deleted{\unsure{$\bigg\}$}}, 
défini ainsi à partir de
\[
 G_0^{+}\subset G_0\subset G
\]
et de l'action de $G$ sur $S$, en termes d'un $G_0^+$-répére
\[
 (s_0,t_0)\in\overrightarrow{A_0}\qquad\text{i.e.}\quad s_0,t_0\in S,\quad\{s_0,t_0\}\in A_0:\tag*{(12)}
\]
\[
\left\{\begin{aligned}
 B'&=G_{s_0}=\{g\in G\mid gs_0=s_0\}\\
 \widetilde U&=B'\cap G_0\deleted{\unsure{{}_{s_0}}}=\{g\in G_0\mid gs_0=s_0\}=\\
 &=\{b'\in B'\mid\delta(b')=1\}\\
 U&=B'\cap G_0^+=\{g\in G_0^+\mid gs_0=s_0\}=\{b'\in B'\mid\tilde\delta(b')=1\}\\
 \widetilde B''&=\{g\in G_0\mid gs_0\in Z.s_0\}\leftarrow\\
 B''&=\widetilde B''\cap G_0^+\quad\deleted{\unsure{\{}}
 \quad[\text{NB }\widetilde B''=\widetilde U.B'']\\
 B&=B'.B''\,(=B'.\widetilde B'')\quad\deleted{\text{NB }\unsure{\{}\,\widetilde B''=\widetilde U.B''}.
\end{aligned}\right.\tag*{(13)}
\]
 \margin{c'est le sous-groupe de $\operatorname{Norm}_{G_0}(\widetilde U)$ contenant $\widetilde U=G_{0s_0}$ qui sert à décrire $Z\simeq\widetilde B''/\widetilde U$ et son action sur $S\simeq G_0/\widetilde U$}
 \margin{\uline{NB} $B$ est bien un \unsure{ss-groupe|sous-groupe|s-groupe} car $B'=G_{s_0}$ est \unsure{stabilisé|normalisé|normalise} par $B''$ qui est défini en termes de $s_0$}
 \margin{produits des sous-ens. dans un groupe $G$\quad $A.B=\{ab\mid a\in A,b=B\}$}}
\scanpage{28}
% archive sheet number: 27
\[
\left.\begin{aligned}
 &\bigl(e(\alpha)u(\lambda)\underbrace{e'(\beta)\bigr)\bigl(e(\alpha')}_{e(\frac{\alpha'}{1+\alpha'\beta})u(\frac{1}{1+\alpha'\beta})e'(\frac{\beta}{1+\alpha'\beta})}u(\lambda')e'(\beta')\bigr)
 \overset{\text{si }\deleted{\alpha'\beta\ne-1}\,1+\alpha'\beta\in\Lambda^*}{=}\\
 &\boxed{=e\left(\alpha+\lambda^2\frac{\alpha'}{1+\alpha'\beta}\right)
 u\left(\frac{\lambda\lambda'}{1+\alpha'\beta}\right)e'\left(\beta'+{\lambda'}^{-2}\frac{\beta}{1+\alpha'\beta}\right)}
\end{aligned}\right\}\quad 1+\alpha'\beta\in\Lambda^*
\]
\[
\left.\begin{aligned}
 &e(\alpha)u(\lambda)\underbrace{e'(\beta)e(\alpha')}_{N(\beta)e'(-\beta)}u(\lambda')e'(\beta')\\
 &\underbrace{\deleted{\unsure{u(\lambda)}}N(\beta)u(\lambda')}_{\in N^-}\underbrace{e'(\beta'-{\lambda'}^{-2}\beta)}_{\in U_1}
 =e(\alpha)\underbrace{u(\lambda)\deleted{\unsure{u}}N(\beta)u(\lambda')}_{N(\beta\lambda^{-1}\lambda')}e'(\beta'-{\lambda'}^{-2}\beta)
\end{aligned}\right\}\quad\begin{gathered}1+\alpha'\beta=0\quad\text{i.e.}\\\alpha'\beta=-1\end{gathered}
\]
\[
e(\alpha)N(\underbrace{\beta\lambda^{-1}\lambda'}_{t})e'(\underbrace{\beta'-{\lambda'}^{-2}\beta}_{\gamma})
=N(t).\underbrace{\underbrace{\deleted{\unsure{e}}N(t)^{-1}\bigl(e(\alpha)\bigr)}_{e'(-t^2\alpha)}e'(\gamma)}_{e'(-t^2\alpha+\gamma)}
\]
\[
e(\alpha)\,t\,e'(\gamma)=\begin{pmatrix}1&\alpha\\0&1\end{pmatrix}
\underbrace{\begin{pmatrix}0&-t^{-1}\\t&0\end{pmatrix}\begin{pmatrix}1&0\\\gamma&1\end{pmatrix}}_{\deleted{\underbrace{\hphantom{00}}\unsure{\bigl(}}\;\begin{pmatrix}-t^{-1}\gamma&-t^{-1}\\t\deleted{\unsure{0}}&0\end{pmatrix}}
=\begin{pmatrix}-t^{-1}\gamma+\alpha t&-t^{-1}\\t&0\end{pmatrix}
\]
\[
\deleted{e(\alpha)N(t)=N(t)e'(\alpha')}
\]
\[
\begin{aligned}
 N(t)\bigl(e(\alpha)\bigr)&\overset{\mathrm{def}}=N(t)e(\alpha)N(t)^{-1}
 =\underbrace{\begin{pmatrix}0&-t^{-1}\\t&0\end{pmatrix}\begin{pmatrix}1&\alpha\\0&1\end{pmatrix}}_{\unsure{=|a}\;\begin{pmatrix}0&-t^{-1}\\t&t\alpha\end{pmatrix}}\begin{pmatrix}0&t^{-1}\\-t&0\end{pmatrix}\\
 &=\begin{pmatrix}1&0\\-t^2\alpha&1\end{pmatrix}\\
 &\boxed{N(t)e(\alpha)=e'(-t^2\alpha)}\qquad\begin{gathered}t\in U_1^{\unsure{+|\ast}}\\\alpha\in U_0\end{gathered}
\end{aligned}
\]
\[
-t^2\alpha+\gamma=\underbrace{-\alpha\beta^2}_{-}\lambda^{-2}{\lambda'}^2+\beta'-{\lambda'}^{-2}\beta
\]
\[
\begin{aligned}
e(\alpha)u(\lambda)e'(\beta)e(\alpha')u(\lambda')e'(\beta')&=\boxed{N(\beta\lambda^{-1}\lambda')e'\bigl(\beta'-{\lambda'}^{-2}\beta-\alpha\beta^2{\lambda'}^2\lambda^{-2}\bigr)}\\
&\text{si }\quad \alpha'\beta=-1\quad\text{i.e.}\quad \alpha'=-\beta^{-1}
\end{aligned}
\]
\scanpage{29}
% archive sheet number: 28
\[
\begin{aligned}
\sigma&\in{B_G}\qquad\deleted{\sigma=}\qquad \sigma t\sigma t\qquad u\\
&\hspace{3em}\sigma=\deleted{t}u.t\qquad u\deleted{\unsure{d}}gu^{-1}=g^{-1}\\
&\hspace{12em}\deleted{ugu^{-1}g^{-1}=g^{-2}}
\end{aligned}
\]
\[
\underbrace{\begin{pmatrix}1&0\\\beta&1\end{pmatrix}\begin{pmatrix}1&-\beta^{-1}\\0&1\end{pmatrix}}_{\begin{pmatrix}1&-\beta^{-1}\\\beta&0\end{pmatrix}}\begin{pmatrix}1&0\\\beta&1\end{pmatrix}
\overset{?}=\begin{pmatrix}0&-\beta^{-1}\\\beta&0\end{pmatrix}\qquad\text{\unsure{ok|OK}}
\]
\[
\begin{aligned}
 \deleted{\unsure{\lambda_1}}u_\lambda\sigma(s_0)&=\lambda^{-1}t_0\\
 \sigma u_\lambda(s_0)&=u_\lambda^{-1}\sigma(s_0)=\lambda t_0
\end{aligned}
\]
% The second hand orientation begins here; a diagonal stroke cancels the passage.
\mspage{\unsure{135|136}}
\deleted{Je \unsure{m'aperçois|m'apperçois} à l'instant que l'image des \inserted{sous-groupes} $T_0$ et de $T_1$ de $T_G$ dans $\operatorname{Aut}H$ \unsure{sont|soit} un seul et \unsure{$\simeq$|m\textsuperscript{e}|$=$} sous-groupe \inserted{et \unsure{égale|égal} à l'image de $T_G$} --- ce qui est évident aussi bien par les formules \deleted{\unsure{qui|que}}
\[
 T_G=T_0.T_H,\qquad \operatorname{Im}\bigl(T_H\to\operatorname{Aut}(H)\bigr)
 \inserted{\subset\operatorname{Im}\bigl(T_0\to\operatorname{Aut}(H)\bigr)}
\]
qui impliquent que $T_0$ et $T_G$ (et \unsure{donc|non} de $\simeq$ $T_1$ et $T_G$ \unsure{\&|?|aussi}, par symétrie) ont $\simeq$ image dans $\operatorname{Aut}(H)$, le sous-groupe en question décidément mérite d'être appellé $\widetilde Z$, donc \uline{\unsure{une utile convention|un allégement|cette convention}} d'écriture !}
\scanpage{30}
% archive sheet number: 29
\mspage{1}
\uline{Structure du groupe} $\mathrm{Sl}(2,\Lambda)$.

1. Soit $M$ un module libre de rang 2 sur un anneau $\Lambda$, avec une structure symplectique
\[
e_M\in(\bigwedge\nolimits^2M)^*.\tag*{(1)}
\]
La donnée d'un tore maximal déployé $T$ de $\mathrm{Sl}(M)$ \uline{revient} à la donnée d'une décomposition de $M$ en somme de deux sous-modules libres \inserted{$L'_i$} de rang 1, formant un ensemble \deleted{$\{L,L'\}$} à deux éléments $\varepsilon_T$ --- \unsure{Il|ce} \deleted{\unsure{leurs|leur}} sont les deux sous-modules propres de $M$, et on a
% slanted
\margin{Si $\operatorname{Spec}\Lambda$ connexe --- Sinon ce n'est le cas que \unsure{lorsqu'on|lorsqu'on a} \uline{\unsure{donne|donné}} un déploiement de $T$, i.e. une paire d'iso $T\simeq G_m$ opposés l'un de l'autre}
\[
\bigotimes_{i\in\varepsilon_T}L'_i\simeq(\bigwedge\nolimits^2M)\wedge_{\pm1}\varepsilon_T\simeq \Lambda\wedge_{\pm1}\varepsilon_T\tag*{(2)}
\]
en utilisant l'isomorphisme
\[
\bigwedge\nolimits^2 M\xrightarrow{\sim}\Lambda\qquad e_M\longmapsto 1\tag*{(3)}
\]
défini par la base $e_M$. L'opération de $T$ sur chaque $L_i$ définit un caractère
\[
T\xrightarrow{\chi_i}\mathbb G_m\qquad i\in\varepsilon_T\tag*{(4)}
\]
qui est un isomorphisme, et (2) implique que l'on a
\[
\prod_{i\in\varepsilon_T}\chi_i=1\tag*{(5)}
\]
i.e. les deux caractères $\chi_i$ sont inverses l'un de l'autre : ce sont les deux isomorphismes de $T$ sur $G_m$ (quand $\operatorname{Spec}\unsure{A|\Lambda}$ est connexe --- sinon on peut dire \inserted{encore} que ce sont les deux iso. ``canoniques'' de $T$ sur $G_m$). Choisir l'un des ces isomorphismes revient donc à choisir
% slanted
\margin{ils sont donnés avec le déploiement}
\scanpage{31}
% archive sheet number: 30
un des deux éléments de $\varepsilon_T$, i.e. à choisir un \uline{ordre} de cet ensemble, qui nous permet dès lors de distinguer \inserted{un} $L_0$ et un $L_1$, de sorte que
\[
 M=L'_0\oplus L'_1,\tag*{(6)}
\]
et qu'on a un isomorphisme d'identification
\[
 T\xrightarrow{\sim}\unsure{\mathbb G|G}_m,\tag*{(7)}
\]
\inserted{un $\lambda\in G_m(\cdot)\simeq$} $T(\cdot)$ opérant sur $L_0$ par l'homothétie $\lambda$, sur $\unsure{L,|L_1}$ par $\lambda^{-1}$.

On peut écrire, à l'aide de (6)
\[
 H\overset{\mathrm{déf}}{=}\underbrace{\mathrm{Sl}(\deleted{\unsure{2,}}M)\unsure{(\Lambda)}}_{=\operatorname{Aut}(M,e_M)}
 \simeq\left\{\begin{pmatrix}a&b\\c&d\end{pmatrix}\;\middle|\;
 \begin{gathered}
 a,d\in\Lambda,\quad b\in\operatorname{Hom}(L'_1,L'_0)\simeq{L'_1}^{-1}L'_0\\
 c\in\operatorname{Hom}(L'_0,L'_1)\simeq{L'_0}^{-1}L'_1,\quad ad-bc=1
 \end{gathered}\right\},\tag*{(8)}
\]
où \deleted{\unsure{$t$}}
\[
 bc=\deleted{\unsure{b}}cb\in\deleted{\unsure{L_1 L_1|\Lambda^*}}\Lambda^* \tag*{(9)}
\]
désigne le composé de $b$ et $c$, identifié à un élément de $\Lambda$ (homothétie par $\lambda$).

La composition des éléments de $\mathrm{\unsure{SL|Sl}}(\unsure{1|2},M)$ est donné par la régle usuelle. Avec ces notations, l'isomorphisme (7) est donné par
\[
 \Lambda^*\xrightarrow{\sim}T(\Lambda)=\left\{\begin{pmatrix}\lambda&0\\0&\lambda^{-1}\end{pmatrix}\;\middle|\;\lambda\in\Lambda^*\right\}\tag*{(10)}
\]
et on pose $u(\lambda)=\begin{pmatrix}\lambda&0\\0&\lambda^{-1}\end{pmatrix}$ pour $\lambda\in\Lambda^*$.

Par la suite, on s'intéresse au groupe des points \deleted{\unsure{$k$}} de $\mathrm{Sl}(M)$ à valeurs \unsure{ds|de} $\Lambda$, plutôt qu'au schéma en groupes, et on notera par la suite par $T$ au lieu de $T(\Lambda)$, le sous-groupe
